\sacramentSummaryStart{Anointing of the Sick}{Second Sacrament of Healing --- Consecration and Comfort in Grave Illness}{Christ the Divine Physician strengthening soul and body and uniting suffering to His Cross}

\begin{summaryidentity}
Definition & Sacrament anointing a gravely ill baptized believer with blessed plant oil and priestly prayer, completing Penance (\emph{ST}~III, q.~65, a.~3; Trent~XIV).\\
Institution & Intimated by Christ in the apostolic mission (Mark~6:12--13) and explicitly promulgated by St.~James (Jas~5:14--15).\\
Sacramental structure & \latin{Sacramentum tantum}: anointing with oil and prayer. \latin{Res et sacramentum}: consecration to Christ crucified. \latin{Res tantum}: comfort, peace, pardon, health if expedient.\\
Order of ends & \textbf{Primary:} spiritual healing, peace, fortitude. \textbf{Secondary:} sin's remnants, pardon if unable to confess, health. \textbf{Contingent:} hope, family comfort. \textbf{Ultimate:} passage to glory.\\
\end{summaryidentity}

\begin{summaryconstitution}
Matter & Olive or other plant oil blessed by bishop (or priest); bodily anointing. & Latin: forehead and hands (c.~1000). East: multiple anointings (\latin{Euchelaion}). & Necessity: single forehead anointing with full form suffices.\\
Form & Deprecatory prayer invoking the Spirit's saving grace (Jas~5:14--15). & Latin (\emph{Sacram Unctionem}): \latin{Per istam sanctam unctionem\ldots\ adiuvet te Dominus gratia Spiritus Sancti, ut a peccatis liberatum te salvet atque propitius allevet}. & Prayer asks the Spirit's grace to free from sin, save, and raise up; Eastern prayers to the Heavenly Physician.\\
Minister & Valid bishop or priest alone (Jas~5:14; Trent~XIV; c.~1003; CCEO~c.~739). & Pastor or priest with care of souls; East classically several priests, validly one. & Deacons and laity never validly anoint; no special faculty needed.\\
Subject & Living baptized person with reason in danger from illness or age (c.~1004). & Given while alert, before grave surgery, or in frailty; repeatable on relapse. & Never for the dead, infants, healthy persons, or obstinate sinners.\\
Change effected & Accidental formal strengthening by grace; relational alignment to the Cross. & \textbf{No character}; repeatable on relapse or graver danger (c.~1004~\S2). & Given in doubt (c.~1005); unconscious if implicitly desired (c.~1006).\\
\end{summaryconstitution}

\begin{summaryeconomy}
Salvation history & Healing oil; Suffering Servant (Is~53:4). & Mark~6:12--13 intimates; Jas~5:14--15 promulgates. & Completes Penance (\emph{ST}~III, q.~65); prepares for Viaticum. & Peaceful passage to beatific glory.\\
Ecclesial order & Consecrates bodily frailty to God's mercy. & Unites personal suffering to Christ's Passion (Col~1:24). & Latin \emph{Anointing of the Sick} (\emph{SC}~73); Eastern \latin{Euchelaion}. & Readies the soul for final glory (\emph{Suppl.}~q.~30).\\
\end{summaryeconomy}

\begin{summarywitnesses}
Scripture & Mark~6:12--13; Jas~5:14--15; Matt~8:17; Is~53:4 & Mark intimates; James promulgates presbyters, oil, prayer, raising, pardon.\\
Fathers & Serapion; Chrysostom; Innocent~I (\emph{Ep.}~25); Bede (\emph{In Marc.}~II.6) & Innocent~I affirms a \latin{genus sacramenti} with episcopal oil; Bede cites Apostles.\\
Scholastics & Aquinas (\emph{ST}~III, q.~65; \emph{Suppl.}~qq.~29--33); Lombard (IV, d.~23) & Orders Anointing to purging sin's debility (\latin{reliquiae peccati}) for glory.\\
Magisterium & Florence~(\S14); Trent~XIV; \emph{Rom.\ Cat.}; \emph{SC}~73; \emph{Sacram Unctionem} & Defines sacrament, restores \emph{Anointing of the Sick}, renews oil and form.\\
Canon law & CIC~1983 (cc.~998--1007); CCEO~1990 (cc.~737--742) & Reserves ministry to priests, sets illness threshold, repeatability, \latin{Euchelaion}.\\
\end{summarywitnesses}

\begin{studybox}{Distinctions to preserve}
\textbf{Anointing vs.\ Viaticum:} Anointing is for \emph{grave illness} (repeatable); Viaticum (Eucharist) is the \emph{sacrament of the dying} (cc.~921--922). \textbf{Priests alone:} only bishops and priests validly anoint (Jas~5:14; c.~1003); deacons and laity cannot. \textbf{Expedient healing:} bodily recovery is granted when willed (\latin{in quantum autem expedit}, Florence~\S14), never replacing union with the Cross.
\end{studybox}

\sacramentSummaryEnd
